% ---
% id: meat-core/frank-wolfe/rem:fo-interface
% kind: rem
% title: "Two reported numbers and one first-order claim"
% summary: >-
%   A run reports its estimate always, and h_lb only where the hypotheses of
%   [[thm:firstorder]] held, membership of the consumed point in the acceptance
%   set times the relaxed set among them. No hypothesis is read of the solver,
%   and withdrawal never weakens what h_lb means.
% prominence: standard
% status: accepted
% uses: [thm:firstorder, def:first-order-setting, def:relaxed-feasible-set, rem:acc-side]
% mentions: [foundations/req:tiers, req:first-order-withdrawal, rem:fo-iterate-checks, rem:fo-returned-point, rem:fo-one-ruler, rem:fo-reachability, rem:fo-check-exposures, rem:fo-oracle-repair, rem:fo-iterate-repair, rem:fo-negative-values, req:verified-feasible-side, rem:fo-arithmetic]
% assumes: []
% refs: []
% provenance: newly-derived
% lean:
% ---

Every run reports its \tref{estimate}, and where the hypotheses of
Theorem~\ref{thm:firstorder} held the \tref{first-order} bound $\hlb$ as
well, naming the \tref{tier} that carries its claim
(\xref{foundations/req:tiers}). The requirements of this submodule say when
$\hlb$ is withdrawn and what a run checks, repairs and records; this remark and
Remarks~\ref{rem:fo-iterate-checks}, \ref{rem:fo-returned-point},
\ref{rem:fo-one-ruler}, \ref{rem:fo-reachability},
\ref{rem:fo-check-exposures}, \ref{rem:fo-oracle-repair},
\ref{rem:fo-iterate-repair} and~\ref{rem:fo-negative-values} give the
mathematics behind them, one topic each.

\begin{remark}\label{rem:fo-interface}
The objective at the point a run returned is an upper estimate by
Theorem~\ref{thm:firstorder}\,(i), respectively~(iiia), which is the part a
run is in: its point is established, where it is established at all, only
in the \tref{relaxed feasible set} (Remark~\ref{rem:fo-iterate-checks}), and
Remark~\ref{rem:fo-returned-point} says what a run records of it, which is
less than the whole premise. Where the hypotheses of that theorem held, the
run reports $\hlb$ as well; Requirement~\ref{req:first-order-withdrawal} says
when it does not.

Among those hypotheses is the membership of the point consumed at an iterate
in the \emph{product} $\acSet\times\Feas^{J}(\varepsilon)$: it is a
hypothesis of Theorem~\ref{thm:firstorder}\,(iiib) as much as (H-POS),
(H-SUPP) and (H-NOPERT). Its two factors are established by different
means: the $J$ factor by the three checks of
Remark~\ref{rem:fo-iterate-checks}, after any repair of
Remark~\ref{rem:fo-iterate-repair}, and the acceptance factor by forming
$\mathbf q_t$ on the \tref{exact island} under ($\alpha$) and ($\beta$) of
Remark~\ref{rem:acc-side}. A failure of either leaves the membership not
established.

No hypothesis is read of the solver. A return that is merely near the
minimum satisfies no inequality in the safe direction, which is why no
bracket pairs the gradient with it: both sides of a bracket are lower
bounds on the two minima over the exact sets, the feasible one verified
from the solver's multipliers in exact arithmetic, the rounding of the
frame its rows are built with deducted
(Requirement~\ref{req:verified-feasible-side}). So a solve that stopped
short of its optimum, or failed, can loosen a bracket but cannot raise
either side above the exact minimum of the linear form at the gradient as
computed, and neither its return nor its status withdraws one; what the
gradient's own binary64 evaluation leaves is
Remark~\ref{rem:fo-arithmetic}'s. Withdrawal is never a weakening of what
$\hlb$ means; there is no reduced form of the claim.
\end{remark}
